%%%PAGE 137 rot=0 legibility=good summary=电磁感应中用动量定理求杆的速度变化/位移/电量的式子及进出磁场速度变化关系；圆周运动离心运动小结
%%%BLOCK id=137-01 ch=12 sec=电磁感应中的动量 kind=derivation alt=none cont=none y=0.04-0.27
% 原稿：几个式子被手绘的大圈/小圈圈住，并有一条斜线从圈引向右上
\[ -B\bar{I}L\,\Delta t=mv_2-mv_1 \]
\[ -BLq=mv_2-mv_1 \]
\[ q=n\,\frac{\Delta\varphi}{\Delta t} \]% CHECK: 分母原稿看似 Δt（常见写法为 R）
\[ -\frac{B^2L^2\bar{v}}{R}\,\Delta t=\Delta p \]
\[ \Delta x=\frac{m\,\Delta v\,R}{B^2L^2} \]
%FIG p137_a 0.54 0.075 0.72 0.16
\notefig[0.25]{figures/p137_a.png}
\[ |v_2-v_1|=|v_3-v_2| \]
\[ q_{\text{进}}=q_{\text{出}} \]
%%%END
%%%BLOCK id=137-02 ch=04 sec=圆周运动·离心运动 kind=knowledge alt=none cont=none y=0.26-0.39
\underline{离转心运}% CHECK: 原稿第2、4字形似“转”“运”，疑为“离心运动”，按原稿转录

$a=\omega^2 r$

$\omega$ 相同 $r$ 越大 $a$ 越大 · $F$\unclear{向}也越大。% CHECK: “向”字原稿形似“方”
%%%END
%%%PAGEEND 137
